\documentclass[12pt, letterpaper]{article}

%-------Packages---------
\usepackage{newpxtext,newpxmath} % Palatino-like text & math (modern)
\usepackage{bboldx}
\usepackage{mathtools}
\usepackage{tikz}
\usetikzlibrary{calc,intersections}
\usepackage{tikz-cd}
\usepackage[all,arc]{xy}
\usepackage{graphicx}

\usepackage{enumerate}
\usepackage[dvipsnames]{xcolor}
\usepackage{float}
\usepackage[left=1in,top=1in,right=1in,bottom=1in]{geometry}
\usepackage[nocheck]{fancyhdr}
\usepackage{multicol}
\usepackage{setspace}

\usepackage{dsfont}
\usepackage{mathrsfs}

\usepackage{tcolorbox}
\tcbuselibrary{theorems,skins,breakable}

\usepackage{xparse}
\usepackage{import}
\usepackage[utf8]{inputenc}

\usepackage{hyperref}
\usepackage{cleveref}

\usepackage{xcolor, ifthen, tcolorbox}
\tcbuselibrary{theorems,skins,breakable}
% Theorem System
% The following environments are provided:
%   New Problem:        \begin{problem} ...
%   New Question Header:   \qh (then followed by subproblems)
%   Subproblem:     \begin{subproblem} ...
%   Problem Proof:  \begin{problemproof} ...
%   Proof:          \\begin{pf}{color} ...
%   Remark:         \rmk (sentence), \rmkb (block)

% Define custom colors
\definecolor{thmscol}{HTML}{F4565B} %For Problems
\definecolor{lemscol}{HTML}{FE844C} %For Lemmes
\definecolor{prpscol}{HTML}{00A7EF} %For proposition
\definecolor{corscol}{HTML}{23DAFF} %For corrolaries
\definecolor{rmkscol}{HTML}{6DEEC5} %For remarks
\definecolor{defscol}{HTML}{18C991} %For definitions
\definecolor{exmscol}{HTML}{7DB800} %For examples
\definecolor{clmscol}{HTML}{165c58} %For claims

%Change between dark(black)/light(white) mode
\newcommand{\colormode}{white}
\ifthenelse{\equal{\colormode}{white}}
      {\newcommand{\TextColor}{black}}
      {\newcommand{\TextColor}{white}}
\newcommand{\pbcolormain}{thmscol!80!\colormode}
\newcommand{\pbcolorsecond}{thmscol!38!\colormode}


% ============================
% Problem Counter
% ============================
\newcounter{subproblemcounter}
\newcounter{problemcounter}

\newcommand{\q}{
    \setcounter{subproblemcounter}{0}%
    \stepcounter{problemcounter} % Increment problem counter
    \color{\TextColor}Problem \arabic{problemcounter}: % Display the problem counter
}

\newcommand{\stepsub}{
    \stepcounter{subproblemcounter}%
    \color{\TextColor}(\alph{subproblemcounter})
}
% ============================
% Problem
% ============================
\newtcolorbox{pb}[1]{
  enhanced,
  frame hidden,
  titlerule=0mm,
  toptitle=1mm,
  bottomtitle=1mm,
  fonttitle=\bfseries\large,
  coltitle=black,
  colbacktitle=\pbcolormain,
  colback=\pbcolorsecond,
  title={\q #1},
  coltext=\TextColor
}

\newtcolorbox{spb}[1]{
  enhanced,
  frame hidden,
  titlerule=0mm,
  toptitle=1mm,
  bottomtitle=1mm,
  fonttitle=\bfseries\large,
  coltitle=black,
  colbacktitle=\pbcolormain,
  colback=\pbcolorsecond,
  title={\stepsub #1},
  coltext=\TextColor,
  attach boxed title to top left={xshift=3mm,yshift=-2mm},
  boxed title style={
    colback=\pbcolormain,
    colframe=\pbcolormain,
  }
}

\newtcolorbox{pbh}[1]{
    enhanced,
    frame hidden,
    titlerule=0mm,
    toptitle=1mm,
    bottomtitle=1mm,
    fonttitle=\bfseries\large,
    coltitle=black,
    colbacktitle=\pbcolormain,
    colback=\pbcolorsecond,
    title={\q #1},
    boxrule=0pt,
    interior hidden,
    attach boxed title to top left={xshift=3mm,yshift=-2mm},
    boxed title style={
        colback=\pbcolormain,
        colframe=\pbcolormain,
    }
}

\newenvironment{problem}[1]{%
  \newpage
  \begin{pb}{#1}
    }{
  \end{pb}
}

\newenvironment{subproblem}[1]{%
  \begin{spb}{#1}
    }{
  \end{spb}
}

\newenvironment{problemproof}{%
    \begin{pf}{\pbcolormain}
        }{
    \end{pf}
}

\newcommand{\qh}[1][]{
    \newpage
    \begin{pbh}{#1}\end{pbh} % Display the problem counter
}

% ============================
% Claim
% ============================

\newtcbtheorem[number within=problemcounter]{claim}{Claim}
{
    enhanced,
    frame hidden,
    titlerule=0mm,
    toptitle=1mm,
    bottomtitle=1mm,
    fonttitle=\bfseries\large,
    coltitle=black,
    colbacktitle=clmscol!50!\colormode,
    colback=clmscol!20!\colormode,
}{clm}

\NewDocumentCommand{\clm}{m+m}{
    \begin{claim*}{#1}{}
        #2
    \end{claim*}
}

\newenvironment{clmpf}{
	{\noindent{\it \textbf{Proof.}}
	\setlength{\parindent}{0pt}
	}
	\tcolorbox[blanker,breakable,left=5mm,parbox=false,
    before upper={\parindent15pt},
    after skip=10pt,
	borderline west={1mm}{0pt}{clmscol!50!\colormode}]
}{
    \textcolor{clmscol!50!\colormode}{\hbox{}\nobreak\hfill$\blacksquare$}
    \endtcolorbox
}

\NewDocumentCommand{\clmp}{m+m+m}{
    \begin{claim*}{#1}{}
        #2
    \end{claim*}

    \begin{clmpf}
        #3
    \end{clmpf}
}


% ============================
% Proof
% ============================

\newenvironment{pf}[1]{%
  \def\pfcolor{#1}% Store the color in a macro
  \noindent{\it \textbf{Proof.}}%
  \tcolorbox[blanker,breakable,left=5mm,parbox=false,%
    before upper={\parindent15pt},%
    after skip=10pt,%
    borderline west={1mm}{0pt}{\pfcolor},%
    coltext=\TextColor
  ]%
  \setlength{\parindent}{0pt}%
}{%
  \textcolor{\pfcolor}{\hbox{}\nobreak\hfill$\blacksquare$}%
  \endtcolorbox%
}

\newtcolorbox{solution}{
  blanker,
  breakable,
  left=5mm,
  parbox=false,%
  before upper={\parindent15pt},%
  after skip=10pt,%
  borderline west={1mm}{0pt}{\pbcolormain},%
  coltext=\TextColor
}


% ============================
% Remark
% ============================


\NewDocumentCommand{\rmk}{+m}{
    {\it \color{rmkscol!100!white}#1}
}

\newenvironment{remark}{
    \par
    \vspace{5pt}
    \begin{minipage}{\textwidth}
        {\par\noindent{\textbf{Remark.}}}
        \tcolorbox[blanker,breakable,left=5mm,
        before skip=10pt,after skip=10pt,
        borderline west={1mm}{0pt}{rmkscol!80!\colormode},
        coltext=\TextColor
        ]
}{
        \endtcolorbox
    \end{minipage}
    \vspace{5pt}
}

\NewDocumentCommand{\rmkb}{+m}{
    \begin{remark}
        #1
    \end{remark}
}


% Define a new command for problems
\renewcommand{\headrule}{\color{\TextColor}\hrulefill}
\newcommand{\C}{\mathbb{C}}
\newcommand{\R}{\mathbb{R}}
\newcommand{\Q}{\mathbb{Q}}
\newcommand{\Z}{\mathbb{Z}}
\newcommand{\N}{\mathbb{N}}
\renewcommand{\P}{\mathbb{P}}
\newcommand{\F}{\mathbb{F}}
\newcommand{\contr}{\Rightarrow \Leftarrow}
\newcommand{\s}{\(\,\)}
\renewcommand{\t}[1]{\text{#1}}

\newcommand{\skcgen}{{\sf Gen}} \newcommand{\skcenc}{{\sf Enc}}
\newcommand{\skcdec}{{\sf Dec}}
\newcommand{\white}{\color{white}} \newcommand{\red}{\color{red}}

% Meta data
\setstretch{1.2}

\begin{document}
\color{\TextColor}
\pagecolor{\colormode}
\setlength{\parindent}{0pt}
\pagestyle{fancy}
\fancyhf{}
\fancyhead[LOE]{\color{\TextColor}\slshape{\textbf{Vincent Ferrara}}}
\fancyhead[ROE]{\color{\TextColor}HW1 --- \thepage}
\fancyhead[C]{\color{\TextColor}EECS 475}

\begin{problem}{}
    Recall that an encryption scheme is formally
  defined by algorithms $\skcgen, \skcenc$, and $\skcdec$
  (respectively, the key generator, encryption, and decryption),
  together with a message space, ciphertext space, and key space.  
\end{problem}
\begin{subproblem}{}
    Give a formal specification for the shift cipher, assuming the message can be any three-letter sequence.
\end{subproblem}
\begin{solution}
    \begin{itemize}
        \item $\mathcal{M}$: $\{0,1,\dots,25\}^3$
        \item $\mathcal{K}$: $\{0,1,\dots,25\}$
        \item $\mathcal{C}$: $\{0,1,\dots,25\}^3$
        \item $\skcgen$: $k \overset{\$}{\leftarrow} \mathcal{K}$
        \item $\skcenc(k,m)$: $c_i = (m_i+k) \mod 26$ where $c_i,m_i$ is the $i$th letter in the message and ciphertext.
        \item $\skcdec(k,m)$: $m_i = c_i-k \mod 26$ where $c_i,m_i$ are the same as defined above.
    \end{itemize}
    Correctness:

    $\skcdec(k,\skcenc(k,m))=m+k-k=m$
\end{solution}

\begin{subproblem}{}
    Give a formal specification for the
      Vigen\`ere cipher\footnotemark. You may assume that the key has a length 5.
    
    \footnotetext{Vigen\`ere cipher can be found in Section 1.1.3 from Vipul Goyal's lecture notes (see Canvas for the link). }
\end{subproblem}
\begin{solution}
    \begin{itemize}
        \item $\mathcal{M}$: $\{0,1,\dots,25\}^5$
        \item $\mathcal{K}$: $\{0,1,\dots,25\}^5$
        \item $\mathcal{C}$: $\{0,1,\dots,25\}^5$
        \item $\skcgen$: $k \overset{\$}{\leftarrow} \mathcal{K}$
        \item $\skcenc(k,m)$: $c_i = (m_i+k_i) \mod 26$ where $c_i,m_i,k_i$ is the $i$th letter in the message, ciphertext, and key.
        \item $\skcdec(k,m)$: $m_i = c_i-k_i \mod 26$ where $c_i,m_i,k_i$ are the same as defined above.
    \end{itemize}

    Correctness:

    $\skcdec(k,\skcenc(k,m))=m+k-k=(m_1+k_1-k_1,\dots,m_n+k_n-k_n)=m$
\end{solution}

\begin{problem}{}
    For each of the following functions, decide whether they are negligible or not in $n$. You only have to provide your answer (whether each one is negligible or not)
and you will get full points if the answer is correct.
You may write your reasoning if you want, and if your answer
is incorrect but your reasoning partly makes sense, 
you may get some partial-credit.
All the logs are two-based.
\end{problem}
\begin{subproblem}{}
    $2^{-\sqrt{n}}$
\end{subproblem}
\begin{solution}
    This is negligible.
\end{solution}

\begin{subproblem}{}
    $(\log n)^{100}$
\end{subproblem}
\begin{solution}
    This is not negligible.
\end{solution}

\begin{subproblem}{}
    $n^{-10}$
\end{subproblem}
\begin{solution}
    This is not negligible.
\end{solution}

\begin{subproblem}{}
    $n^{-\log n}=(2^{\log n})^{- \log n}= 2^{-(\log n)^2}$
\end{subproblem}
\begin{solution}
    This is negligible.
\end{solution}

\begin{subproblem}{}
    $2^{-\log n \cdot \log( \log n)}$
\end{subproblem}
\begin{solution}
    This is negligible.
\end{solution}

\begin{problem}{}
    Prove that if a private-key encryption scheme $(\skcgen,\skcenc,\skcdec)$ satisfies Shannon secrecy, then it satisfies perfect secrecy.
\end{problem}
\begin{problemproof}
    Assume a private-key encryption scheme satisfies Shannon secrecy.

    Let $m \in \mathcal{M}$ and $c \in \mathcal{C}$, and let $M,C$ be random variables for the messages, ciphertext respectively.

    Consider
    \[\P[C = c \mid M = m] = \frac{\P[C=c \t{ and } M=m]}{\P[M=m]}=\frac{\P[M=m \mid C=c]\P[C=c]}{\P[M=m]}\]
    From Shannon secrecy we have
    \[=\frac{\P[M=m]\P[C=c]}{\P[M=m]}=\P[C=c]\]
    Since we did this for a general message $m$ we can also do this process with a $m' \in \mathcal{M}$ which implies
    \[=\P[C = c \mid M= m']\]
    Therefore, the private-key encryption scheme satisfies perfect secrecy.
\end{problemproof}

\qh
\begin{subproblem}{}
    Suppose that for some private-key encryption scheme,
the ciphertext always starts with the letter ``c'', 
can this be a perfectly secure encryption?
Please explain why.
\end{subproblem}
Yes this can be perfectly secure.

\begin{problemproof}
    Assume $\Sigma=(\skcgen, \skcenc, \skcenc)$ is a perfectly secure encryption.

    Define a new scheme $\Sigma'=(\skcgen', \skcenc', \skcenc')$ s.t. the message space, and the key space is the same, but the ciphertext space is the same except add a $c$ to the start of all the strings.

    Let $\skcgen' = \skcgen$, $\skcenc'(k,m)= c + \skcenc(k,m)$ where $+$ means append, and $\skcdec'(k,c+y)=\skcdec(k,y)$.
    
    This is correct since for $k$ output by $\skcgen'$, and $m \in \mathcal{M}$ we have

    \[\skcdec'(k, \skcenc'(k,m))=\skcdec'(k,c+\skcenc(k,m))=\skcdec(k,\skcenc(k,m))=m\]
    
    Let $m,m' \in \mathcal{M}$, and $x \in \mathcal{C}$ for $\Sigma'$. By definition of $\Sigma'$ this implies that $x=c+y$ for some $y$ in the ciphertext space of $\Sigma$.

    Consider
    \begin{align*}
        \P[\skcenc'(K,m)=x] &= \P[\skcenc(K,m)=y]
        \intertext{This is by definition of what $\skcenc'$ is.
        Also since $\Sigma$ is perfectly secure this implies that}
        \P[\skcenc(K,m)=y] &= \P[\skcenc(K,m')=y] = \P[\skcenc'(K,m')=x]
    \end{align*}
    Therefore, $\Sigma'$ is perfectly secure, and has all ciphertext's start with c.
\end{problemproof}

\begin{subproblem}{}
    Suppose that for some private-key encryption scheme,
the key always starts with the letter ``c''; 
can this be a perfectly secure encryption?
Please explain why.
\end{subproblem}
Yes this can be perfectly secure.

\begin{problemproof}
    Assume $\Sigma=(\skcgen, \skcenc, \skcenc)$ is a perfectly secure encryption.

    Define a new scheme $\Sigma'=(\skcgen', \skcenc', \skcenc')$ s.t. the message space, and the ciphertext space is the same, but the key space is the same except add a $c$ to the start of all the strings.

    Let $\skcgen' = c + \skcgen$, $\skcenc'(c+k,m)=\skcenc(k,m)$, and $\skcdec'(c+k,c)=\skcdec(k,c)$.
    
    This is correct since for $c+k$ output by $\skcgen'$, and $m \in \mathcal{M}$ we have

    \[\skcdec'(c+k, \skcenc'(c+k,m))=\skcdec(k,\skcenc(k,m))=m\]

    Let $m,m' \in \mathcal{M}$, and $x \in \mathcal{C}$ for $\Sigma'$.

    Consider
    \begin{align*}
        \P[\skcenc'(K'=c+K,m)=x] &= \P[\skcenc(K,m)=x]
        \intertext{This is by definition of what $\skcenc'$ is.
        Also since $\Sigma$ is perfectly secure this implies that}
        \P[\skcenc(K,m)=x] &= \P[\skcenc(K,m')=x] = \P[\skcenc'(c+K=K',m')=x]
    \end{align*}
    Therefore, $\Sigma'$ is perfectly secure, and has all ciphertext's start with c.
\end{problemproof}

\begin{subproblem}{}
    Suppose a message is a string containing letters from 
``a'' to ``z''. If the ciphertext always contains the first letter of the message, can it be perfectly secure encryption? 
Please explain why.
\end{subproblem}
Yes this can be perfectly secure.

\begin{problemproof}
    Assume $\Sigma'(\skcgen, \skcenc, \skcenc)$ is a perfectly secure encryption.

    Let $H="abcdefghijklmnopqrstuvxyz"$

    Define a new scheme $\Sigma'=(\skcgen', \skcenc', \skcenc')$ s.t. the message space, and the key space is the same, but the ciphertext space is the same except add $H$ to the start of all the strings.

    Let $\skcgen'=\skcgen$, $\skcenc'=H+\skcenc$, and $\skcdec'(k,H+y)=\skcdec(k,y)$.

    The proof for why this is correct and perfectly secure is the same as in (a) just replace the $c$ with $H$.

    Also, this is a scheme where the first letter is always in the ciphertext since $H$ contains all the letters $"a" to "z"$.
\end{problemproof}

\begin{problem}{}
    Let $\mathbb{Z}_n$ denote integers modulo $n$,
and let $\mathbb{Z}_n \setminus \{0\}$ denote all the elements of $\mathbb{Z}_n$ except $0$ (e.g. $\mathbb{Z}_3 = \{0, 1, 2\}$ and $\mathbb{Z}_3 \setminus \{0\} = \{1, 2\}$). 
% Definition
For some fixed message length $\ell$,
let key space $\mathcal{K}$, message space $\mathcal{M}$, and ciphertext space $\mathcal{C}$ all be $(\mathbb{Z}_{n} \setminus \{0\})^\ell$ (i.e. vectors of length $\ell$ over elements from $\mathbb{Z}_{n} \setminus \{0\}$).
We partially define the following encryption scheme over $\mathbb{Z}_n \setminus \{0\}$:
\begin{align*}
    \mathsf{Gen}(1^\ell) &: (k_1, \dots, k_\ell) \overset{\$}{\leftarrow} (\mathbb{Z}_{n} \setminus \{0\})^\ell\\
    \mathsf{Enc}(k, m) &: (k_1 \cdot m_1, k_2 \cdot m_2, \dots, k_\ell \cdot m_\ell)\mod n
\end{align*}
\end{problem}

\begin{subproblem}{}
    Define a decryption algorithm for the above encryption scheme over $\mathbb{Z}_3 \setminus \{0\}$
    (i.e. $\mathcal{K} = \mathcal{M} = \mathcal{C} = \mathbb{Z}_{3} \setminus \{0\}$), 
    and prove its correctness.
\end{subproblem}
\begin{solution}
    Since $1*1=1$ and $2*2=1 \implies 2^{-1}=2$ the multiplicative inverse this implies that every non-zero element of $\Z_3$ has a multiplicative inverse.

    \noindent Therefore, $k_i^{-1}$ is well-defined.

    \noindent Let $\skcdec(k,c): (k_1^{-1} \cdot c_1,\dots, k_\ell^{-1} \cdot c_{\ell}) \mod 3$

    \noindent Correctness:

    \noindent Let $k$ be in the output of $\skcgen(1^\ell)$, and $m \in \mathcal{M}$.
    \[\skcdec(k, \skcenc(k,m))=\skcdec(k, c)\]

    \noindent Since shown above every non-zero element has an inverse this implies that $m_i = k_i^{-1} \cdot c_i$ Thus,
    \[=\skcdec(k, c)=(k_1^{-1}\cdot c_1,\dots, k_\ell^{-1} \cdot c_\ell)=m\]
\end{solution}

\begin{subproblem}{}
    Show that the encryption scheme over $\mathbb{Z}_{3} \setminus \{0\}$ is perfectly secret.
\end{subproblem}
\begin{problemproof}
    Let $m,m' \in \mathcal{M}$, $c \in \mathcal{C}$.

    Let $m=(m_1,\dots,m_\ell)$. Since each $m_i$ is invertible this implies that there is a unique key

    $k = c \odot m^{-1}$ where $\odot$ is component wise multiplication. This implies that $k \odot m = c$ Therefore,

    \begin{align*}
        \P[\skcenc(K,m)=c]=\P[K=k]=(1/2)^\ell
    \end{align*}

    Since $m$ is general this also holds for  $m'$. Thus,

    \[\P[\skcenc(K,m)=c]=\P[\skcenc(K,m')=c]\]
\end{problemproof}

\begin{subproblem}{}
    Suppose we define the encryption scheme over $\mathbb{Z}_4 \setminus \{0\}=\{1,2,3\}$. 
    Is this a valid encryption scheme?
\end{subproblem}
\begin{problemproof}
    This is not a valid encryption scheme.

    Consider if $k_i=2$ and $c_i=2$. Since $k_i \cdot 1 = 2 = k_i \cdot 3$.

    This leads to use not being able to invert the message from the ciphertext since in this situation under the encryption scheme both $m_i=1$ and $m_i=3$ leads to this ciphertext.

    Thus, we cannot fully decrypt the ciphertext given the key which implies it is not a valid encryption.
\end{problemproof}

\begin{problem}{}
    Prove or give a counter example: An encryption scheme with a
  finite message space $\mathcal{M}$ is perfectly secure if and only
  if for every probability distribution $\mathcal{D}$ over
  $\mathcal{M}$ and every $c_0, c_1 \in \mathcal{C}$ over the cipher
  space $\mathcal{C}$, we have $\Pr[C = c_0] = \Pr[C = c_1]$, where
  $C$ is the ciphertext corresponding to a message chosen according to
  $\mathcal{D}$.
\end{problem}
\begin{problemproof}
    Let $\Sigma=(\skcgen, \skcenc, \skcdec)$ be an encryption scheme s.t. $\mathcal{M}$ is finite, and $\mathcal{C}$ is the set of ciphertexts that can be returned by the encryption function.

    Assume for every distribution $\mathcal{D}$ over $\mathcal{M}$ and all $c,c' \in \mathcal{C}$,
    \[\P[C=c]=\P[C=c']\]
    
    Since $\mathcal{M}$ is finite this implies that $\mathcal{C}$ is also finite since we only take this space to be the output of $\skcenc$.

    Therefore, $\sum_{c \in \mathcal{C}} \P[C=c]=1 \implies |\mathcal{C}| \P[C=c] = 1 \implies \P[C=c]=\dfrac{1}{|\mathcal{C}|}$

    Thus, for any message distribution it will give $\mathcal{C}$ the uniform distribution.

    This implies that $\P[C=c \mid M=m]=\P[C=c]$ since $C,M$ are independent since the distribution over $\mathcal{M}$ has no effect on the distribution over $\mathcal{C}$.

    Therefore, for $m,m' \in \mathcal{M}$ and $c \in \mathcal{C}$

    \[\P[C=c \mid M=m] = \P[C=c] = \P[C=c \mid M=m']\]

    Thus, we have perfect secrecy.

    The converse to this statement is not true.

    Let $\mathcal{M} = \{0,1\}$, $\mathcal{C}=\{a,b,c\}$, and $\mathcal{K}=\{k_1,\dots,k_6\}$.

    Define $\skcgen$ by uniformly picking one of the keys.

    Define $\skcenc$ by
    \begin{center}
        \begin{tabular}{c|cc}
            $k$ & $\skcenc(0)$ & $\skcenc(1)$\\
            \hline
            $k_1$  & $a$ & $b$\\
            $k_2$  & $a$ & $b$\\
            $k_3$  & $a$ & $c$\\
            $k_4$  & $b$ & $a$\\
            $k_5$  & $b$ & $a$\\
            $k_6$  & $c$ & $a$\\ 
        \end{tabular}
    \end{center}

    $\P[C=a]=1/2$, $\P[C=b]=1/3$, and $\P[C=c]=1/6$.

    $\P[C=a \mid M = 0] = \P[K=k_1 \t{ or } K=k_2 \t{ or } K=k_3]=1/6+1/6+1/6=1/2$

    $\P[C=a \mid M = 1] = 1/2$

    $\P[C=b \mid M = 0] = 1/3 = \P[C=b \mid M = 1]$

    $\P[C=c \mid M = 0] = 1/6 = \P[C=c \mid M = 1]$

    Therefore, this encryption scheme has perfect secrecy, but $\P[C=a] \neq \P[C=b]$.
\end{problemproof}

\begin{problem}{}
    Consider the probability ensembles $\{X_n\}_n$ and 
$\{Y_n\}_n$. Consider distinguishers $\mathcal{D}_1, \ldots, \mathcal{D}_m$ where 
$m = \t{poly}(n)$ 
for some polynomial function $\t{poly}(\cdot)$, 
such that only one can distinguish $\{X_n\}_n$ and $\{Y_n\}_n$ 
with non-negligible probability, but you do not know which one. 
Construct a distinguisher $\mathcal{D}^*$ that can distinguish between $\{X_n\}_n$ 
and $\{Y_n\}_n$ with non-negligible probability. 
Please describe your construction of $\mathcal{D}^*$ and provide a proof. 
\end{problem}
\begin{problemproof}
    Construct $\mathcal{D}^*$ s.t. pick $i \leftarrow \{1,\dots, m(n)\}$ uniformly at random independent of the input.

    Then output $D_i$

    Let $\Delta_i(n)=\P[x \leftarrow X_n \mid D_i(x)=1]-\P[y \leftarrow Y_n \mid D_i(y)=1]$.

    By assumption there exists a unique index $j$ and a polynomial $c \in \R_{>0}$ with $\mid \Delta_j(n) \mid \geq 1/n^c$,
    while for all $i \neq j$, $|\Delta_i| \leq \t{negl}(n)$.

    Thus, 
    \[|\P[x \leftarrow X_n \mid D^*(x)=1]-\P[y \leftarrow Y_n \mid D^*(y)=1]|= \left\lvert 1/m \sum_{i=1}^m \Delta_i(n)\right\rvert \geq 1/m(|\Delta_j(n)| - |\Sigma_{i \neq j} \Delta_i(n)|) \]

    Since a polynomial times a negligible function is negligible this implies that 
    
    $(m-1) \t{negl}(n)$ is negligible. Thus, for sufficiently large $n$,

    \[(m-1) \t{negl}(n) \leq \dfrac{1}{2n^c}\]

    This implies that,
    \[|\P[x \leftarrow X_n \mid D^*(x)=1]-\P[y \leftarrow Y_n \mid D^*(y)=1]| \geq (1/m)(1/n^c - 1/2n^c) = \frac{1}{2m n^c}\]

    Since $m$ and $n^c$ are polynomial this implies that $m \cdot n^c$ is a polynomial. Thus, $\frac{1}{2m \cdot n^c}$ is non-negligible. Thus, $D^*$ distinguishes $\{X_n\}_n$ from $\{Y_n\}_n$ with non-negligible probability.
\end{problemproof}
\end{document}